\section{Implementation Details}
\label{appendix:implementation-details}


\subsection{Experimental Configuration}
\label{appendix:experiment-config}
To ensure the reproducibility of our experiments, all tests were conducted with consistent parameter configurations. The random seed for the large language model was fixed to \textbf{123}, and the temperature parameter was set to \textbf{0} to eliminate randomness in the generation process. Additionally, GPT-4o is employed as our LLM-based evaluator. Furthermore, all dense embedding computations employed the "\textbf{nomic-ai/nomic-embed-text-v2-moe}" model \cite{Chunk4} as the foundational embedding model. For tokenization operations, we uniformly adopted the "\textbf{nomic-ai/nomic-embed-text-v2-moe}" as tokenizer, with all token-based chunks standardized to \textbf{750} tokens in length. For detailed parameter settings of other experiments, please refer to \Cref{appendix:parameter-configuration}.

\subsection{Compute Resources}
\label{appendix:compute-resources}
In our experiments, we employed \textbf{GPT-4o-mini} \cite{gpt-4o-mini} and \textbf{Deepseek-V3} \cite{deepseek-v3} as the underlying large language models (LLMs). (All operations utilizing large language models (GPT-4o, GPT-4o-mini, Deepseek-V3) are completed by invoking APIs). Notably, the proposed experiments can be readily replicated on ordinary machines, ensuring broad accessibility and reproducibility for the research community. Experiments were executed on a dedicated research workstation configured with an AMD Ryzen 7 7800X3D 8-Core Processor, an NVIDIA GeForce RTX 4070 Ti SUPER GPU, and \(64\) GB of DDR5 RAM. 

\subsection{Evaluation Metrics}
Our efficiency is quantified in our paper through token consumption. Our decision not to use and can not use precise wall-clock time latency as a direct metric stems from two key reasons: (1) the strong correlation between LLM inference time and the number of output tokens, and (2) the inherent uncertainties and fluctuations caused by network latency during API calls to LLMs, which would introduce inaccuracies into direct time measurements. 


\subsection{Parameter Configuration}
\label{appendix:parameter-configuration}

The system involves the following key parameters:

\begin{itemize}[leftmargin = 15pt]
    \item \textbf{Node Correlation Weight ($\alpha$)}: 
    This parameter adjusts the system's reliance on edge embedding for determining strong links. Experimental validation suggests setting $\alpha$ within the range of 0.1--0.2. In this study, we adopt $\alpha = \bm{0.1}$.

    \item \textbf{Strong Connection Threshold ($\lambda$)}: 
    As discussed in \Cref{sec:theoretical-analysis}, we typically set it to a value below 0.775, but this is only under theoretical conditions. In practice, various factors should be comprehensively considered—setting it too low may lead to increased retrieval costs due to a lower memory threshold outweighing the noise reduction, while setting it too high may result in reduced memory capacity. Through empirical testing, we consider values between 0.5 and 0.75 to be reasonable. In the experiments of this paper, we select this value as \textbf{0.55}.
\end{itemize}

Other critical parameters include:

\begin{itemize}[leftmargin = 15pt]
    \item \textbf{Synonym Similarity Threshold}: 
    Used during knowledge graph construction to merge entities whose embedding similarity exceeds this threshold. This value should be adjusted according to the specific embedding model characteristics. Our experiments employ \textbf{0.7} as the default value.

    \item \textbf{Maximum Hop Count}: 
    Controls the maximum number of nodes (hops) during subgraph expansion for word queries. We set this parameter to \textbf{10} to balance computational efficiency and information completeness. For other baselines with similar configurations, we uniformly apply the same setting for fair comparison.

    \item \textbf{Question Decomposition Limit}: 
    Determines the maximum number of sub-questions for semantic decomposition of user queries. To maintain comparability with baseline methods, we set this to \textbf{1} (i.e., preserving the original question without decomposition).

    \item \textbf{Initial Seed Node Count}: 
    Specifies the number of seed nodes in the query initialization phase. We configure this as \textbf{2} seed nodes (selecting the two most query-relevant nodes), consistent with relevant baseline studies.
\end{itemize}

\subsection{Implementation of Baselines}
\textbf{Traditional retrieval methods:}
\begin{itemize}[leftmargin=*]
\item \textbf{BM25} \cite{BM25}. A classic information retrieval algorithm improving on TF-IDF \cite{TF-IDF}, it estimates document-query relevance via a probabilistic framework. It effectively handles long documents and short queries by integrating term frequency (TF), inverse document frequency (IDF), and document length normalization to calculate relevance scores, enhancing retrieval accuracy through adaptive term weighting. Following its principles, we have realized the algorithm.

\item \textbf{NaiveRAG} \cite{Base2}. The NaiveRAG paradigm, an early "Retrieve-Read" framework, involves indexing multi-format data into vector-encoded chunks stored in a database, retrieving top-K similar chunks via query vector similarity, and generating responses by prompting large language models with queries and selected chunks. We have developed the corresponding code based on its principles.


\end{itemize}

\textbf{Traditional KG-RAG systems}:
\begin{itemize}[leftmargin=*]
\item \textbf{HippoRAG2} \cite{Main6}. This framework enhances retrieval-augmented generation (RAG) to mimic human long-term memory, addressing factual memory, sense-making, and associativity. It uses offline indexing to build a knowledge graph (KG) with triples extracted by an LLM, integrating phrase nodes (concepts) and passage nodes (contexts) via synonym detection. Online retrieval links queries to KG triples using embeddings, filters irrelevant triples with an LLM (recognition memory), and applies Personalized PageRank (PPR) on the KG for context-aware passage retrieval. Improvements include dense-sparse integration, deeper query-to-triple contextualization, and effective triple filtering.
\item \textbf{LightRAG} \cite{Main2}. This method addresses limitations of existing RAG systems (e.g., flat data representations, insufficient contextual awareness) by integrating graph structures into text indexing and retrieval. It employs a dual-level retrieval paradigm: low-level retrieval for specific entities/relationships and high-level retrieval for broader themes, combining graph structures with vector representations to enable efficient keyword matching and capture high-order relatedness. The framework first segments documents, uses LLMs to extract entities/relationships, and constructs a knowledge graph via LLM profiling and deduplication for efficient indexing. 
\item \textbf{GraphRAG} \cite{Main3}. This approach aims to enable large language models (LLMs) to perform query-focused summarization (QFS) over large private text corpora, especially for global sensemaking queries that require understanding of the entire dataset. It constructs a knowledge graph in two stages: first, an LLM extracts entities, relationships, and claims from source documents to form a graph where nodes represent key entities and edges represent their relationships. Second, community detection algorithms partition the graph into a hierarchical structure of closely related entity communities. The LLM then generates bottom-up community summaries, with higher-level summaries recursively incorporating lower-level ones. Given a query, GraphRAG uses a map-reduce process: community summaries generate partial answers in parallel (map step), which are then combined and summarized into a final global answer (reduce step). 
\end{itemize}
For these three baselines, we use their default hyperparameters (except for chunk size, which is uniformly set to 750, consistent with \ours) and prompts.

\textbf{LLM-guided KG-RAG system}:
\begin{itemize}[leftmargin=*]
\item \textbf{Plan-on-Graph (PoG)} \cite{KG1}. This paradigm aims to address the limitations of existing KG-augmented LLMs, such as fixed exploration breadth, irreversible paths, and forgotten conditions. To achieve this, it decomposes the question into sub-objectives and repeats the process of adaptive reasoning path exploration, memory updating, and reflection on self-correcting erroneous paths until the answer is obtained. PoG designs three key mechanisms: 1) Guidance: Decomposes questions into sub-objectives with conditions to guide flexible exploration. 2) Memory: Records subgraphs, reasoning paths, and sub-objective statuses to provide historical information for reflection. 3) Reflection: Employs LLMs to decide whether to self-correct paths and which entities to backtrack to based on memory.We adapt their code from their GitHub repository\footnote{\url{https://github.com/liyichen-cly/PoG}} to use the KG constructed by \ours~(More details can be seen at \Cref{appendix:::some-details-of-PoG}).

\end{itemize}

For all methods described above, the "nomic-ai/nomic-embed-text-v2-moe" model is uniformly used as the embedding model. Methods incorporating knowledge graph memory functionality utilize ChromaDB for knowledge graph storage.

\subsection{Implementation of LLM-as-judgment}
\label{appendix::answer-evaluation}

Followed previous work \cite{Benchmark1},We employ the prompt shown in \Cref{tab:check-ans-prompt} to conduct LLM-as-judge (GPT-4o) based answer evaluation.


\subsection{Additional Experiments on More Baseline Models}
\label{appendix::More Baseline Models}
For better understanding, we consider an additional baselines MemoRAG \cite{Main11}. This framework enhances RAG to tackle long-context processing, overcoming conventional RAG’s reliance on explicit queries and well-structured knowledge. Inspired by human cognition, it uses a dual-system architecture: a light long-range system builds global memory (via configurable KV compression) optimized by RLGF (rewarding clues that boost answer quality), and an expressive system generates final answers. For tasks, the light system produces draft clues to guide retrieving relevant evidence from long contexts, then the expressive system outputs the final answer. It outperforms advanced RAG methods (e.g., GraphRAG \cite{Main3}) and long-context LLMs in both QA and non-QA tasks, with balanced inference efficiency.   


\textbf{Experimental Setup}. Given MemoRAG \cite{Main11}'s capability to manage memory for ultra-long texts, we focused our evaluation on the \textit{LongDependencyQA} task of the LooGLE dataset \cite{Benchmark1} in this experiment; for the baseline configuration, the model employed is \textit{gpt-4o-mini}, the embedding\textunderscore model is uniformly set to \textit{nomic-ai/nomic-embed-text-v2-moe} (consistent with our experimental setup), and all other parameters adhere to MemoRAG's default settings—including the mem\textunderscore mode parameter, which adopts its default value of \textit{memorag-mistral-7b-inst}. The pre-trained weights of \textit{memorag-mistral-7b-inst} were obtained from the Hugging Face Hub, and we verified the model's license agreement (Apache 2.0) on its official Hugging Face Model Card to ensure compliance with research usage requirements, avoiding potential issues related to intellectual property or usage restrictions.


\begin{table}[htbp]
  \centering
  \captionsetup{skip=10pt}           % Increase space between table and caption (adjust "15pt" as needed)
  \begin{tabular}{l|l|c}
    \hline
    \hline
    \textbf{Methods } & \textbf{Baseline Model} & \textbf{Long-Dependency-Task (\%)} \\
    \hline
    MemoRAG        & gpt-4o-mini             & 46.97                              \\
    HippoRAG2      & gpt-4o-mini             & 39.60                              \\
    REMINDRAG      & gpt-4o-mini             & 57.04                              \\
    \hline
    \hline
  \end{tabular}
  \caption{Performance on Long-Dependency-Task (Accuracy $\uparrow$).}
  \label{tab:more-baseline-experiment}
\end{table}

\textbf{Experimental Results}. As shown in Table \ref{tab:more-baseline-experiment}, MemoRAG indeed demonstrates strong performance in long-dependency tasks: its 46.97\% accuracy on the Long-Dependency-Task outperforms HippoRAG2's 39.60\% by 7.37 percentage points, a result supported by its dual-system architecture that efficiently builds global memory and optimizes clue selection via RLGF. However, it still lags significantly behind our \ours~ method—our approach achieves 57.04\% accuracy, a 10.07-percentage-point lead over MemoRAG, as MemoRAG's KV compression may lose fine-grained long-range information, while \ours~ better captures dynamic, long-distance associations critical to the task.